
% ============================================================
% PROYECTO DE TRABAJO DE INVESTIGACIÓN - UTP
% Plantilla LaTeX para Overleaf
%
% Para cambiar el estilo de citado:
%   APA : \documentclass[citestyle=apa]{utp-plan}
%   IEEE: \documentclass[citestyle=ieee]{utp-plan}
%
% Bibliografía: BibTeX + archivo referencias.bib
% En Overleaf: Menu > Settings > Compiler: pdfLaTeX
% ============================================================

\documentclass[citestyle=ieee]{utp-plan}


\usepackage{color}
\usepackage{soul}


% ---------------- Datos de la carátula ----------------
% Reemplace el texto de ejemplo y conserve los comandos.
\campus{UTP Chiclayo}
\carrera{Ingeniería de Sistemas e Informática}
\titulo{Optimización de arquitectura de Deep Learning para el reconocimiento de conductas anómalas: Detección eficiente de shoplifting en Edge Computing para Retail}
\integranteuno{Huanca Fernández, Kevin (U22244540)}
\integrantedos{Llacsahuanga Buques, Diego Alexis (U22203352)}
\docente{Alfaro Purisaca, Paul Anthony}
\ciudad{Chiclayo, Perú}
\ciclo{Agosto II 2026}

\begin{document}
	
	% chat: 20260913-ZD783LDV
	
	\UTPportada
	\UTPindice
	
	% ============================================================
	% 1. PROBLEMA DE INVESTIGACIÓN
	% ============================================================
	\section{Problema de investigación (máximo 1/2 página)}
	
En los últimos años, las pequeñas y medianas empresas como minimarkets y bodegas enfrentan un severo problema de robo en tiendas, que causa pérdidas de inventario y merma económica en el sector retail peruano. Gran parte de la seguridad con la que cuentan dichos establecimientos depende únicamente de sistemas de videovigilancia tradicionales (CCTV), los cuales presentan varias limitaciones operativas; entre ellas, están la fatiga visual de los operadores humanos, la presencia de puntos ciegos y su naturaleza aleatoria. Al mismo tiempo, han surgido diversas técnicas de aprendizaje profundo (Deep Learning) como una alternativa novedosa y prometedora; sin embargo, la literatura científica revela tres vacíos que es importante resaltar:

\begin{enumerate}[leftmargin=*, noitemsep, topsep=2pt] 
		\item \textbf{Falta de modelado de conductas pre-delictivas:} La mayoría de los modelos se centran unicamente en detectar el momento exacto en donde ocurre el robo de forma reactiva, ignorando lo más importante, que seria los patrones de movimiento, los cuales caracterizan y son de gran ayuda para identificar la intención delictiva de la persona.
\item \textbf{Elevado costo computacional:} Las arquitecturas actuales requieren de hardware costoso, como servidores y GPUs de alto rendimiento, ocasionando que los comercios minoristas locales no puedan adquirir dicho hardware debido a su presupuesto limitado.
\item \textbf{Escasez de investigación en el contexto nacional:} Existe una gran escasez de estudios y benchmarks adaptados a la infraestructura, iluminación y ángulos de cámara del retail peruano, identificándose unicamente el antecedente de \cite{aguero_intelligent_2025}.
	\end{enumerate}
	
	
	% ============================================================
	% 2. OBJETIVO GENERAL Y ESPECÍFICOS
	% ============================================================
	\section{Objetivo general y específicos (máximo 1 página)}
	
	\subsection{Objetivo general}
	Desarrollar un modelo de aprendizaje profundo basado en la arquitectura YOLO + LSTM y el framework PCB (\textit{Pre Crime Behavior}), optimizado para bajo costo computacional y desplegable en edge devices, para mejorar las métricas de rendimiento en la detección temprana de conductas de \textit{shoplifting} en el entorno retail minorista peruano.
	
	\subsection{Objetivos específicos}

		\begin{enumerate}
			\item \textbf{Construir y estructurar un dataset propio} con diversas secuencias de video de comportamiento de compra y tambien conductas sospechosas en diferente comercios minoristas locales, aplicando diferntes técnicas de preprocesamiento, segmentación de secuencias y data augmentation.
			
			\item \textbf{Diseñar e integrar la arquitectura híbrida YOLO + LSTM}, combinando redes CNN y LSTM para la extracción de características espaciales y el modelado temporal de la secuencia pre-delictiva, consiguiendo tambien optimizar hiperparámetros como la longitud de secuencia y el tamaño de lote.
			
			\item \textbf{Implementar y evaluar el despliegue del modelo sobre edge diveces} de bajo consumo computacional como en tarjetas NVIDIA Jetson o Raspberry Pi, analizando las diferente métricas de eficiencia operativa como la latencia de detección y los fotogramas por seguno (FPS).
			
			\item \textbf{Validar la precisión y rendimiento predictivo del sistema} mediante las diversas métricas estándar como Precisión, Recall, F1-Score y mAP, comparando los resultados obtenidos frente a los detectores espaciales estáticos y arquitecturas pesadas reportadas en la literatura.
		\end{enumerate}
	
	% ============================================================
	% 3. JUSTIFICACIÓN
	% ============================================================
	\section{Justificación}
	
La presente investigación se justifica por la necesidad de los establecimientos pequeños del sector retail peruano frente a las conductas asociadas al shoplifting, buscando desarrollar una alternativa eficiente para detectar rápidamente la presencia de hurtos en tiendas. En la actualidad, la mayoría de negocios dependen de soluciones de vigilancia convencionales, los cuales sufren de algunas limitaciones como: la fatiga visual, puntos ciegos y la dificultad para detectar patrones de comportamientos previos al hurto. A esto se suma que las soluciones basadas en aprendizaje profundo suelen requerir hardware costoso, lo que dificulta su adopción en este tipo de comercios.

Con respecto a la problemática mencionada, se propone una arquitectura de Deep Learning basado en la integración de la arquitectura YOLO con LSTM y el framework PCB (Pre Crime Behavior), enfocado a reconocer patrones espacio temporales en las conductas previas al crimen. La implementación estará enfocada a dispositivos de bajo consumo como NVIDIA Jetson o Raspberry Pi, buscando facilitar su aplicación en los comercios minoristas con bajo presupuesto. Esta propuesta se encuentra alineada con el objetivo general de desarrollar un modelo optimizado para mejorar la detección temprana de shoplifting en el contexto del retail minorista peruano.

Asimismo, la investigación busca aportar evidencia adaptada al contexto peruano por medio de la construcción de un dataset propio que incluya diversas secuencias de comportamientos sospechosos en diferentes condiciones de comercios minoristas locales. Esto permitirá considerar factores importantes como las variaciones de iluminación, ángulos de cámara, oclusiones y puntos ciegos.

Los resultados esperados ayudarán a determinar si la arquitectura propuesta puede lograr alcanzar un equilibrio entre la capacidad detección y la eficiencia computacional, proporcionando información sobre su viabilidad en la implementación para entornos de retail minorista peruano.
	
	% ============================================================
	% 4. REVISIÓN DE LA LITERATURA ACTUAL O ESTADO DEL ARTE
	% ============================================================
	\section{Estado del arte}
  	
	\subsection{Introducción}
  
  La detección automática de shoplifting en centros comerciales mediante técnicas de \textit{computer vision} y deep learning ha experimenado crecimiento constante en los últimos años, principalmente impulsado debido a las crecientes pérdidas económicas que asociadas al hurto minorista así como por los avances en arquitecturas de redes neuronales convolucionales y recurrentes. La revisión de literatura evidencia una evolución: pasamos desde los primeros acercamientos basados en \textit{hand-crafted features} como HOG (\textit{Historigram of Gradients}) y LBP (\textit{Local Binary Pattern}) (\cite{m_deep_2023}) combinadas con clasificadores clásicos, pasando a modelos espacio temporales profundos como 3DCNN, LRCN, y CNN-LSTM, hasta los detectores de una sola etapa como la arquitectura YOLO, capaces de operar en hardware limitado. Resaltamos que al revisar de manera crítica los diez trabajos más relevantes en este campo, se observa que la mayoría de estudios utilizan el dataset UCF-Crime, lo que limita la aplicabilidad y generalización de estos al contexto peruano.
  
  
	\subsection{Desarrollo}
	
	\subsubsection{Modelos Convolucionales de Una Sola Etapa Orientados a Tiempo Real}
  
    Uno de los ejes de la literatura más prevalente es el uso de arquitecturas YOLO (\textit{You Only Look Once}) para la detección de comportamiento
    anómalo o sospechoso. En 2025 Hameed et al desarrolló un sistema basado en YOLOv8 entrenado con el dataset UCF-Crime más un conjunto de
    imágenes recolectadas por los autores, se alcanzó un mAP50-95 del 94.07\% y una precisión (\textit{Precision}) del 94.93\%, manteniendo el tiempo de respuesta en tan solo     12 milisegundoss por fotograma, esto demuestra que este tipo de arquitecturas son viables para ser implementadas en una solución de tiempo real en edge devices \cite{hameed_real_2025}. En otro trabajo realizado por Pavithra et al en 2024 (\cite{v_deep_2024}) compara modelos espacio-temporales (LRCN, ConvLSTM e InceptionV3) contra YOLOv8 encontrando que este obtuvo un mejor desempeño en la métrica de mAP50 de 83.2\% justamente por prescindir del análisis temporal, lo que lo hace menos costoso computacionalmente.
    
    Otro trabajo que continúa esta línea de investigación es el precedente geográfico más cercano al presente estudio, desarrollado por Aguero y Rengifo en 2025 (\cite{aguero_intelligent_2025}), plantea un sistema de
    videovigilancia inteligente para minimercados del distrito de Santiago de Surco, Lima, empleando YOLOv8 sobre un dataset propio de más de 2,700
    imágenes reales y simuladas, con resultados de 80\% en Precision y Recall, sin embargo se encuentra una caída del mAP50 del 84\% frente al mAP50-95 del 57\%, es esto lo que deja en evidencia la baja generalización del modelo frente a condiciones más difíciles (poca iluminación, variaciones de la cámara y oclusión).
    
    Si bien, ambos trabajos confirman que YOLO es la arquitectura más aplicable para el despliegue en tiempo real en un entorno con un hardware prácticamente limitado, es importante hacer notar que el único antecedente peruano \cite{aguero_intelligent_2025} limita su dataset a un solo distrito de Lima y a un tamaño reducido de 2782 imágenes, a comparación de la mayoría de estándares internacionales, que suelen tener más de 15000 imágenes \cite{hameed_real_2025}. Además, frente a la caída de la métrica mAP50-95 con respecto a mAP50, se sugiere que el modelo no es completamente generalizable en condiciones de variabilidad lumínica, lo cual abre una oportunidad concreta para este trabajo: ampliar y diversificar el dataset con datos provenientes de distintos comercios minoristas peruanos podría mejorar sustancialmente la robustez del modelo sin incrementar necesariamente el costo computacional, dado que YOLO ya ha demostrado ser eficiente en este aspecto.
		
	\subsubsection{Modelos Híbridos Espacio-Temporales Orientados a la Detección Temprana}
  
    Un segundo grupo de investigaciones se centra en capturar la dimensión temporal del comportamiento sospechoso, partiendo de la premisa de que
    el hurto no es un evento instantáneo sino un proceso que se desarrolla a lo largo de varios segundos, incluso minutos, antes de consumarse.
    Una de las soluciones con este enfoque combina una red LSTM con una arquitectura InceptionV3, optimizada mediante el
    ajuste de la longitud de secuencia y el tamaño de lote, logrando un 89.36\% de precisión en un dataset sintetizado por ellos mismos, aunque a costa de
    tiempos de entrenamiento que, en sus configuraciones más profundas, superaron las cuatro horas por experimento (\cite{ansari_optimized_2023}). Por su parte,
    Martínez-Mascorro et al., propone un enfoque distinto: en lugar de detectar el hurto en sí, modelan el comportamiento previo
    al delito (pre-crime behavior) mediante redes 3DCNN, obteniendo una precisión de hasta 80.76\% en el dataset existente hasta y 89\% en
    uno sintetizado por los mismos autores, documentando de forma detallada el impacto de la resolución de imagen y la profundidad
    temporal (número de frames) sobre el tiempo de entrenamiento y la exactitud del modelo (\cite{martinez-mascorro_suspicious_2021}).
    
    Dado que uno de nuestros objetivos específicos es \textit{implementar y evaluar el despliegue del modelo sobre edge devices}, los resultados de esta sección nos muestran que, a mayor profundidad temporal y resolución, mejora la
    exactitud pero se dispara el tiempo de entrenamiento y, por extensión, el costo de inferencia en producción. Para un contexto de comercio
    minorista peruano, estas soluciones resultan poco viables tal como están planteadas en la literatura. Esto refuerza la necesidad de optar
    por una arquitectura ligera (como YOLO) antes que por modelos híbridos profundos, y de evaluar explícitamente el trade-off entre exactitud
    y eficiencia computacional como parte de la metodología de la tesis, algo que ningún estudio revisado aborda de forma sistemática en conjunto
    con un dataset local.
    
  \subsubsection{Enfoques orientados a la privacidad y mitigación de sesgos}
    Un tercer eje, más reciente, responde a las crecientes preocupaciones éticas sobre el uso de video crudo en sistemas de vigilancia.   
    Ante esto, se introduce el dataset \textit{PoseLift}, un dataset construido en colaboración cercana con una tienda retail minorista de Estados     Unidos que, en lugar de     almacenar video, conserva únicamente secuencias de pose humana, anonimizando así la identidad de los clientes, 
    los autores evalúan modelos de detección de anomalías no supervisados (GEPC, TSGAD, STG-NF), obteniendo como máximo un AUC-ROC de 67.46\% lo que es
    inferior a las cifras de exactitud reportadas por los enfoques supervisados de detección de objetos. Los propios autores reconocen que este desempeño
    más modesto se debe, en parte, al desbalance de clases inherente al problema y a la dificultad de discriminar acciones ambiguas, como el hecho de sacar
    un teléfono del bolsillo que puede confundir al modelo con ocultar un objeto \cite{rashvand_exploring_2025}.
    
    Se evidencia entonces que la preservación de la privacidad y la detección de hurtos no
    necesariamente son incompatibles, asimismo que existe una brecha de desempeño entre los enfoques que preservan la privacidad y los que priorizan
    el rendimiento bruto. Para efectos de nuestra tesis, este antecedente sugiere que, si bien la privacidad es un aspecto deseable a mediano plazo,
    priorizar en esta primera etapa un enfoque supervisado basado en detección de objetos (más maduro y con métricas más altas y estables) resulta más
    pertinente para el objetivo de mejorar las métricas de rendimiento planteado, dejando el análisis de anonimización de datos como una línea de trabajo
    futura una vez validado el dataset peruano.
  
  \subsubsection{Sistemas híbridos con integración de sensores e IoT}
  
    Finalmente, un grupo de trabajos complementa la visión por computadora con sensórica adicional para reforzar la robustez de la detección.
    En \cite{m_deep_2023} diseñaron un sistema para joyerías que combina YOLO con sensores de peso (celda de carga) y sensores infrarrojos de
    proximidad, validado sobre un prototipo con Raspberry Pi 3, logrando una exactitud de validación de 97.73\%, aunque sobre un dataset muy reducido, de
    tan solo 144 imágenes. Varun et al \cite{m_intelligent_2025}, proponen un sistema similar orientado a bancos y joyerías, que combina detección YOLO
    durante el día con un módulo de detección de sonido basado en MFCC durante la noche, e integra actuadores IoT (ESP32) para activar alarmas
    y liberar gas no tóxico ante una amenaza detectada.
    
    Hay una limitación metodológica clara para
    los fines de esta tesis en los artículos de esta sección: los datasets utilizados son extremadamente pequeños (144 imágenes en el caso de \cite{m_deep_2023}) y las métricas
    de validación no siempre corresponden a condiciones representativas de un comercio real, sino a prototipos de laboratorio o tiendas
    simuladas. Esto confirma, una vez más, que la escasez de datasets reales y de tamaño adecuado sigue siendo el cuello de botella principal
    en esta línea de investigación, independientemente del país o del enfoque tecnológico elegido.


  \subsection{Conclusiones del Estado del Arte}
  
  Se encontraron diversos puntos centrales en los estudios. Los diez trabajos revisados coinciden en un
  diagnóstico común, el hurto en tiendas es un problema que crece sostenidamente en términos económicos (\cite{rashvand_exploring_2025}, \cite{v_deep_2024}, \cite{ansari_optimized_2023}) y que los sistemas de videovigilancia tradicionales, dependientes del monitoreo humano, resultan insuficientes debido a que dependen completamente de la atención del operador la cual se ve superada completamente por la cantidad de datos que debe procesar (\cite{anikwe-onyiwa_recent_2025}). A partir de ahí, existe también coincidencia en que el aprendizaje profundo, ya sea mediante detectores de una sola etapa como YOLO o mediante arquitecturas espacio temporales (3DCNN, LRCN, CNN-LSTM) constituyen la vía más prometedora para automatizar esta detección, y en que el dataset UCF-Crime es, hasta la fecha, el recurso más utilizado para entrenar y comparar estos modelos (\cite{martinez-mascorro_suspicious_2021}, \cite{chinthulla_deep_2025}, \cite{hameed_real_2025}).
  
Sin embargo, los estudios difieren notablemente en la forma de plantear el problema. Mientras un grupo de investigaciones \cite{hameed_real_2025, v_deep_2024, aguero_intelligent_2025} trata la detección como un problema de clasificación binaria en tiempo real sobre imágenes o fotogramas individuales, otro grupo \cite{martinez-mascorro_suspicious_2021, ansari_optimized_2023, chinthulla_deep_2025} insiste en la necesidad de modelar la dimensión temporal del comportamiento, es decir, detectar la intención antes de que el crimen se consume, lo que exige arquitecturas más costosas de entrenar. Esta discrepancia en la literatura no ha sido resuelta: no existe consenso sobre si conviene priorizar la velocidad de inferencia (enfoque espacial) o la riqueza temporal (enfoque espacio-temporal), y los propios resultados numéricos entre ambos enfoques no son directamente comparables porque se evalúan sobre datasets, métricas y condiciones de prueba distintas \cite{rashvand_exploring_2025}.

Se aprecia una tendencia clara hacia la adopción de YOLO en sus versiones más recientes (v5, v8, v11) como arquitectura preferida cuando el objetivo es el despliegue en tiempo real \cite{hameed_real_2025, aguero_intelligent_2025, m_intelligent_2025, m_deep_2023}, desplazando gradualmente a los modelos puramente recurrentes en los estudios más nuevos. También se observa una tendencia creciente, aunque aún incipiente, hacia la preocupación ética: los trabajos más recientes \cite{hameed_real_2025, rashvand_exploring_2025} discuten explícitamente el riesgo de sesgo demográfico y de falsos positivos, algo ausente en los estudios de 2020-2022.

La ausencia más significativa es la escasez de datasets construidos con data real de tiendas, capturada en campo y en condiciones no controladas. De los diez estudios, solo dos \cite{rashvand_exploring_2025} en Estados Unidos y \cite{aguero_intelligent_2025} en Lima trabajaron con grabaciones reales de un comercio; el resto recurre a UCF-Crime (videos agregados de internet, sin contexto de tienda específico) o a escenarios escenificados por actores en laboratorio \cite{chinthulla_deep_2025, m_deep_2023, m_intelligent_2025}. Para el contexto peruano, la ausencia es aún más marcada: existe un único antecedente \cite{aguero_intelligent_2025}, limitado a un solo distrito y a un dataset de tamaño moderado.

Es importante destacar que estas limitaciones no son una interpretación nuestra, sino que los propios autores las reconocen explícitamente: el modelo de cuatro clases de Chinthulla \cite{chinthulla_deep_2025} obtuvo un F1-score de 0.00 en la clase "Theft" por la escasez de muestras disponibles, mientras que el mejor modelo de Rashvand et al. (2025) confunde acciones cotidianas, como sacar el celular del bolsillo, con intentos de ocultamiento, lo que confirma que el desbalance de clases continúa siendo un obstáculo no resuelto en ambos estudios. De forma similar, el mAP50-95 de 57\% reportado por Solórzano y Candela \cite{aguero_intelligent_2025} refleja una degradación del desempeño ante oclusión y condiciones de iluminación variables, por lo que los propios autores sugieren ampliar el dataset y explorar arquitecturas más recientes. Ansari y Singh \cite{ansari_optimized_2023} y Martínez-Mascorro et al. \cite{martinez-mascorro_suspicious_2021}, por su parte, coinciden en que una mayor profundidad temporal eleva considerablemente el tiempo de entrenamiento, sin que ello garantice una mejora proporcional en la exactitud.

A pesar de los avances documentados, no existe en el contexto peruano un dataset de hurto minorista que combine tres condiciones a la vez: haber sido recolectado directamente en comercios reales (bodegas, minimarkets y tiendas de conveniencia) de más de una localidad; tener un volumen suficiente para sostener un entrenamiento robusto sin el desbalance de clases que ha limitado a estudios previos; y haber sido evaluado explícitamente bajo el criterio de costo computacional, y no solo de exactitud.

El único antecedente local \cite{aguero_intelligent_2025} evidencia una caída de desempeño bajo condiciones reales de oclusión y variabilidad de iluminación (mAP50-95 de 57\%), lo que sugiere que ampliar y diversificar la data de entrenamiento con muestras propias del comercio minorista peruano podría mejorar la robustez del modelo sin necesidad de recurrir a arquitecturas más pesadas.

  Considerando estos vacíos, se plantea la siguiente pregunta de investigación: ¿En qué medida el diseño de un modelo de deep learning basado en YOLO + LSTM y el framework PCB permite mejorar las métricas de rendimiento en la detección temprana de shoplifting (precisión, recall, mAP50-95) manteniendo un bajo costo computacional ejecutable en edge devices, dentro del entorno retail minorista peruano?

		
	% ============================================================
	% 5. MARCO TEÓRICO
% ============================================================
\section{Marco Teórico}

El desarrollo de soluciones automatizadas para la prevención del hurto en el comercio minorista requiere un marco conceptual integrado que transite, de manera lógica y jerárquica, desde los fundamentos de la videovigilancia computacional hasta las arquitecturas algorítmicas de aprendizaje profundo y sus métricas de evaluación operacional en dispositivos de cómputo limitado.

\subsection{Videovigilancia y Seguridad en el Entorno Retail}

\subsubsection{Videovigilancia Inteligente (\textit{Intelligent Video Surveillance})}

Los sistemas tradicionales de Circuito Cerrado de Televisión (CCTV) han operado históricamente como mecanismos pasivos de registro, cuya eficacia depende críticamente de la supervisión humana en tiempo real. Sin embargo, la fatiga visual, la saturación de monitores y la limitación de la atención selectiva reducen drásticamente la capacidad operativa de los guardias de seguridad para detectar eventos anómalos en el momento en que ocurren.

Frente a estas deficiencias, la Videovigilancia Inteligente (\textit{Intelligent Video Surveillance} -- IVS) se constituye como el paradigma orientado a transformar los flujos de video continuos en información estructurada mediante el procesamiento analítico autónomo \cite{anikwe-onyiwa_recent_2025}.

La función primordial de un sistema IVS reside en la automatización del monitoreo mediante algoritmos capaces de interpretar escenas complejas sin intervención humana constante, transitando de una seguridad puramente reactiva, donde la grabación se consulta únicamente tras la consumación de una pérdida, hacia una prevención proactiva en tiempo real.

En el sector \textit{retail}, la implementación de la IVS busca mitigar la merma desconocida (\textit{inventory shrinkage}), la cual engloba las pérdidas de inventario derivadas del hurto sin violencia. La característica distintiva de las soluciones modernas de IVS es su capacidad para operar de forma continua, filtrando el ruido visual de las actividades rutinarias de los clientes y aislando patrones de movimiento relevantes para la seguridad del establecimiento \cite{hameed_real_2025}.

\subsubsection{Reconocimiento de Actividades Humanas (\textit{Human Activity Recognition})}

Dentro del dominio de la visión por computadora, el Reconocimiento de Actividades Humanas (\textit{Human Activity Recognition} -- HAR) constituye la disciplina encargada de identificar, clasificar e interpretar las acciones ejecutadas por uno o más individuos a partir de secuencias de imágenes o flujos de video.

Mientras que la detección de objetos se restringe a responder qué elementos están presentes en la escena y dónde se ubican, el HAR aborda la dimensión dinámica, respondiendo a la interrogante de qué acción está realizando el sujeto a lo largo del tiempo \cite{ansari_optimized_2023}.

Las actividades humanas se categorizan comúnmente según su nivel de abstracción temporal en tres niveles: gestos (movimientos atómicos de extremidades), acciones (secuencias simples como caminar, inclinarse o tomar un producto) y actividades complejas (cadenas de acciones interrelacionadas dispuestas en un contexto temporal).

En el contexto del comercio minorista, el HAR enfrenta el desafío de discriminar entre interacciones comerciales legítimas, tales como examinar la etiqueta de un producto, comparar artículos o depositarlos en una canasta de compra, e interacciones anómalas o potencialmente delictivas, como el ocultamiento rápido de un objeto en prendas de vestir o bolsas personales.

La consecución de esta tarea requiere articular descriptores espaciales que capten la postura y la interacción mano-objeto con modelos temporales que analicen la progresión de dicha conducta en el tiempo \cite{chinthulla_deep_2025}.


\subsection{Modelado Conductual Pre-Delictivo}

\subsubsection{El Marco del Comportamiento Pre-Delictivo (\textit{Pre-Crime Behavior})}

La detección convencional de delitos mediante aprendizaje profundo ha tendido a enfocarse en la clasificación del acto ilícito en el instante preciso de su consumación.

Sin embargo, en el ámbito del hurto minorista, señalar el delito únicamente cuando el producto ya ha sido ocultado reduce el margen de acción del personal de seguridad, transformando la alerta automatizada en una confirmación tardía de la pérdida.

Para resolver esta falencia, Martínez-Mascorro et al. \cite{martinez-mascorro_suspicious_2021} formalizaron el marco del Comportamiento Pre-Delictivo (\textit{Pre-Crime Behavior} -- PCB), un enfoque conceptual que reinterpreta el análisis de video al modelar la secuencia de conductas preliminares que preceden a la ejecución del robo.

El marco PCB postula que la comisión de un hurto no es un evento instantáneo e aislado, sino el resultado de un proceso conductual estructurado en etapas temporales bien definidas. Dichas fases comprenden:

\begin{enumerate}
    \item Momento de Primera Aparición (\textit{First Appearance Moment} -- FAM): representa el ingreso y desplazamiento natural del individuo en el espacio comercial.

    \item Momento Comprensivo del Crimen (\textit{Comprehensive Crime Moment} -- CCM): constituye la etapa pre-delictiva propiamente dicha, caracterizada por acciones sutiles de preparación o exploración, tales como la verificación constante del entorno, giros de cabeza repetitivos para ubicar cámaras o personal, y la manipulación inusual de productos o prendas.

    \item Momento Estricto del Crimen (\textit{Strict Crime Moment} -- SCM): corresponde al acto físico del hurto o apoderamiento del objeto sin intención de pago.

    \item Momento Posterior al Hurto (\textit{Post-Theft Moment} -- PTM) y Retorno a la Normalidad (\textit{Back to Normality Moment} -- B2NM): fases en las cuales el sujeto busca disimular la acción y reincorporarse al flujo ordinario de clientes para abandonar el establecimiento.
\end{enumerate}

La ventaja teórica y práctica de adoptar el marco PCB radica en desplazar la frontera de detección hacia la fase CCM.

Al entrenar modelos de aprendizaje profundo capaces de reconocer los patrones de movimiento característicos del comportamiento sospechoso previo al robo, el sistema de videovigilancia genera alertas tempranas que permiten la intervención disuasiva antes de que la mercancía sea sustraída \cite{chinthulla_deep_2025}.


\subsection{Arquitecturas de Aprendizaje Profundo Espacio-Temporales}

\subsubsection{Extracción de Características Espaciales}

El procesamiento analítico de flujos de video requiere, en primera instancia, transformar la información pictórica no estructurada de cada fotograma en representaciones vectoriales densas que condensen la semántica de la escena.

Las Redes Neuronales Convolucionales (\textit{Convolutional Neural Networks} -- CNN) representan el pilar fundamental para esta tarea. A través de operaciones de convolución, agrupación (\textit{pooling}) y funciones de activación no lineales, las CNN extraen de forma jerárquica características visuales de bajo nivel (bordes, texturas y gradientes de color) hasta abstraer patrones de alto nivel, tales como siluetas humanas, regiones de interés o categorías de productos \cite{anikwe-onyiwa_recent_2025}.

En tareas que exigen la localización precisa de sujetos en tiempo real, los detectores de una sola etapa (\textit{single-stage detectors}) de la familia YOLO (\textit{You Only Look Once}) han reemplazado a las arquitecturas tradicionales de dos etapas debido a su alta eficiencia computacional.

Al dividir la imagen de entrada en una grilla regular y predecir simultáneamente los recuadros delimitadores (\textit{bounding boxes}) y las probabilidades de clase, algoritmos como YOLOv8 logran clasificar y localizar personas u objetos a altas tasas de fotogramas por segundo \cite{hameed_real_2025}.

Sin embargo, como señalan Pavithra et al. \cite{v_deep_2024} y Solórzano y Candela \cite{aguero_intelligent_2025}, los detectores espaciales puros operan sobre fotogramas de manera aislada, careciendo de la dimensión de memoria necesaria para distinguir si un movimiento identificado en una imagen fija constituye un gesto de compra legítimo o el inicio de una maniobra de ocultamiento.


\subsubsection{Modelado de Secuencias Temporales y Memoria Recurrente}

Dado que las conductas analizadas bajo el marco PCB se despliegan a lo largo del tiempo, la extracción de características espaciales debe complementarse con modelos capaces de procesar secuencias temporales.

Las Redes Neuronales Recurrentes (\textit{Recurrent Neural Networks} -- RNN) introducen conexiones con realimentación que permiten mantener un estado oculto interno, actuando como una memoria del contexto pasado.

No obstante, las RNN convencionales sufren del problema de desvanecimiento o explosión del gradiente al procesar secuencias prolongadas, lo que limita su capacidad para aprender dependencias a largo plazo.

Para superar esta restricción, Hochreiter y Schmidhuber introdujeron las redes de Memoria a Corto y Largo Plazo (\textit{Long Short-Term Memory} -- LSTM).

Una unidad LSTM incorpora una celda de memoria regulada por tres mecanismos de compuerta:

\begin{itemize}
    \item La compuerta de olvido (\textit{forget gate}), que decide qué información previa descartar.

    \item La compuerta de entrada (\textit{input gate}), que selecciona qué nuevos datos almacenar en el estado de la celda.

    \item La compuerta de salida (\textit{output gate}), que determina el valor emitido hacia el siguiente paso temporal.
\end{itemize}

La integración funcional entre la extracción espacial convolucional y el modelado temporal recurrente se consolida en la arquitectura de Redes Convolucionales Recurrentes de Largo Plazo (\textit{Long-Term Recurrent Convolutional Network} -- LRCN).

En una red LRCN, una serie de fotogramas continuos es procesada individualmente por un módulo CNN (\textit{TimeDistributed}), el cual genera una secuencia de vectores de características. Esta secuencia es posteriormente canalizada hacia capas LSTM que evalúan la evolución temporal de la conducta.

Como demostró Chinthulla \cite{chinthulla_deep_2025}, la arquitectura LRCN resulta especialmente apta para el análisis del hurto comercial, puesto que la componente convolucional captura la interacción visual entre el cliente y el entorno, mientras que la componente LSTM procesa la transición suave entre las etapas del comportamiento pre-delictivo.


\subsection{Computación en el Borde y Eficiencia Computacional}

\subsubsection{Procesamiento en Dispositivos de Borde (\textit{Edge Computing})}

El paradigma tradicional de analítica de video ha dependido de la transmisión de flujos de video continuos hacia servidores centralizados en la nube o estaciones de trabajo de alto rendimiento equipadas con GPUs dedicadas.

No obstante, en el contexto del comercio minorista de pequeña y mediana escala (minimarkets y bodegas), esta arquitectura centralizada resulta inviable debido a los costos recurrentes de ancho de banda, la latencia de red en la transmisión y la falta de infraestructura de cómputo avanzada \cite{aguero_intelligent_2025}.

La Computación en el Borde (\textit{Edge Computing}) plantea descentralizar el procesamiento, ejecutando los algoritmos de inferencia directamente en dispositivos de hardware locales integrados al punto de captura de video.

Al procesar la información en tarjetas de cómputo embebido de bajo consumo, tales como las plataformas NVIDIA Jetson o Raspberry Pi, se elimina la dependencia de la nube, se reduce la latencia de respuesta a milisegundos y se resguardan las garantías de privacidad al evitar la transferencia externa de imágenes de videovigilancia \cite{hameed_real_2025, pranav_deep_2023}.


\subsubsection{Costo Computacional y Métricas de Inferencia}

El despliegue de modelos de aprendizaje profundo en dispositivos de borde impone restricciones estrictas respecto a la complejidad algorítmica y el consumo de recursos.

El costo computacional de un modelo se evalúa formalmente mediante indicadores de eficiencia tanto estructurales como de inferencia en tiempo de ejecución:

\begin{enumerate}
    \item Volumen de parámetros: determina la cantidad de pesos entrenables de la red y su huella de memoria RAM/VRAM en el dispositivo local.

    \item Latencia de inferencia: representa el tiempo requerido (medido en milisegundos) por el sistema para procesar una entrada y emitir un diagnóstico de clasificación.

    \item Tasa de procesamiento (\textit{Frames Per Second} -- FPS): indica el número de fotogramas por segundo que el modelo es capaz de evaluar continuamente. Para considerarse operativo en tiempo real en sistemas CCTV, la tasa de procesamiento debe situarse en el rango de 15 a 30 FPS.
\end{enumerate}

Como argumentan Ansari y Singh \cite{ansari_optimized_2023}, la optimización del costo computacional en modelos espacio-temporales no solo depende de la ligereza de la red base, sino del ajuste preciso de hiperparámetros como la longitud de la secuencia (\textit{Sequence Length}) y el tamaño de lote (\textit{batch size}), los cuales dictan la ventana de memoria necesaria para evaluar la conducta sin saturar la capacidad de cómputo local.


\subsection{Métricas de Evaluación de Rendimiento}

\subsubsection{Métricas de Clasificación y Detección Espacio-Temporal}

La validación empírica de un sistema de detección de hurtos exige métricas matemáticas rigurosas que contemplen tanto la corrección en la clasificación conductual como la precisión en la localización espacial.

En problemas de clasificación binaria o multiclase donde las muestras de delitos son sustancialmente menores que los eventos de compra normales, la exactitud simple (\textit{Accuracy}) resulta insuficiente por verse distorsionada por el desbalance de clases \cite{rashvand_exploring_2025}.

Por ello, la evaluación se fundamenta en las siguientes métricas derivadas de la matriz de confusión:

\begin{enumerate}
    \item Precisión (\textit{Precision}): mide la proporción de alertas de sospecha correctamente identificadas respecto al total de alertas emitidas por el modelo, siendo el indicador crítico para evaluar la minimización de falsos positivos:

    \begin{equation}
        \text{Precisión} = \frac{TP}{TP + FP}
    \end{equation}

    donde $TP$ representa los verdaderos positivos y $FP$ los falsos positivos.

    \item Exhaustividad o Sensibilidad (\textit{Recall}): calcula la proporción de eventos delictivos reales detectados por el sistema sobre el total de hurtos efectivamente acontecidos, midiendo la capacidad para mitigar los falsos negativos ($FN$):

    \begin{equation}
        \text{Recall} = \frac{TP}{TP + FN}
    \end{equation}

    \item Puntuación F1 (\textit{F1-Score}): constituye la media armónica entre la Precisión y el Recall, ofreciendo una métrica de desempeño global balanceada en escenarios desproporcionados:

    \begin{equation}
        \text{F1-Score} =
        2 \times
        \frac{\text{Precisión} \times \text{Recall}}
        {\text{Precisión} + \text{Recall}}
    \end{equation}

    \item Precisión Media Promedio (\textit{Mean Average Precision} -- mAP): en modelos que incorporan detección de objetos o delimitación de sujetos, la mAP evalúa el área bajo la curva de precisión-recall en función del solapamiento espacial (\textit{Intersection over Union} -- IoU) entre el recuadro predicho y la anotación de referencia.

    Se evalúa principalmente mediante mAP50 (con un umbral IoU de 0.50) y mAP50-95 (promediado en umbrales de IoU desde 0.50 hasta 0.95 en incrementos de 0.05), siendo esta última la prueba más estricta sobre la solidez del encuadre visual bajo situaciones de oclusión parcial \cite{aguero_intelligent_2025}.
\end{enumerate}
		
	% ============================================================
	% 6. HIPÓTESIS
	% ============================================================
	\section{Hipótesis}
	% El diseño de un modelo de deep learning basado en la arquitectura híbrida
  % YOLO + LSTM + framework PCB, entrenado con un dataset propio de comercios minoristas
  % peruanos, permite superar el mAP50-95 de 57\% reportado por el antecedente local
  % \cite{aguero_intelligent_2025}, manteniendo una latencia de inferencia adecuada para
  % su ejecución en edge devices (Jetson Nano/Raspberry Pi).
  
	% ============================================================
	% 7. METODOLOGÍA
	% ============================================================
	\section{Metodología}
	
	% ============================================================
	% 8. CRONOGRAMA DE TRABAJO
	% ============================================================
	\section{Cronograma de trabajo}
	
	% ============================================================
	% 9. PRESUPUESTO
	% ============================================================
	\section{Presupuesto}
		
	% Espacio para que el alumno agregue su presupuesto.
	
	% ============================================================
	% 10. BIBLIOGRAFÍA
	% ============================================================
	\section{Bibliografía}
	
	\bibliography{referencias}
	
	% ============================================================
	% 11. ANEXOS
	% ============================================================
	\section{Anexos}
		
	% Para anexos reales puede utilizar:
	% \appendix
	% \section{Nombre del anexo}
	% Contenido del anexo...
	
\end{document}